% VI30 detailed challenge solutions -- VI/24+ proof-depth standard.

% VI30-DETAILED-CHALLENGE-01
\subsection*{Challenge VI.30.C1 solution}
\begin{solution}
The statement
\[
\dim U=\dim X
\qquad
(\varnothing\ne U\subseteq X\text{ open})
\]
is not a consequence of any one of the most obvious topological adjectives.

\textbf{Irreducible is not enough.}
Take
\[
X=\Spec k[x,y]_{(x)}.
\]
This scheme is integral and therefore irreducible. It is Noetherian and has dimension \(1\). Yet
\[
D(x)
\]
contains only the generic point, hence has dimension \(0\).

\textbf{Equidimensionality of components is not enough.}
The same example is integral, so it has only one irreducible component and is equidimensional in the
componentwise sense, but the dimension still drops on \(D(x)\).

\textbf{Jacobson alone is not enough.}
The scheme
\[
X=\Spec k\sqcup\mathbb A^1_k
\]
is finite type over \(k\), hence Jacobson, but the open-and-closed component
\[
U=\Spec k
\]
has dimension \(0\) while \(\dim X=1\).

\textbf{Finite type over a field alone is not enough.}
The same reducible example is finite type over \(k\), yet open-dimension equality fails.

A clean sufficient hypothesis is precisely the variety hypothesis used in VI/30:
\[
X\text{ integral and finite type over }k.
\]
Then \(X\) is Noetherian; every nonempty open \(U\) is quasi-compact, integral, and finite type over \(k\), hence a
variety. Moreover
\[
K(U)=K(X).
\]
Therefore
\[
\dim U
=
\operatorname{trdeg}_kK(U)
=
\operatorname{trdeg}_kK(X)
=
\dim X.
\]

The examples show why the theorem should not be remembered merely as an ``irreducible-space'' fact. It is a
finite-type algebraic-geometric theorem whose proof is controlled by the common function field and the
dimension--transcendence-degree formula.
\end{solution}

% VI30-DETAILED-CHALLENGE-02
\subsection*{Challenge VI.30.C2 solution}
\begin{solution}
Dominance gives an inclusion
\[
K(Y)\hookrightarrow K(X).
\]
Because \(X\) and \(Y\) are varieties, both function fields are finitely generated over \(k\). For a tower of
finitely generated field extensions
\[
k\subseteq K(Y)\subseteq K(X),
\]
transcendence degrees satisfy
\[
\operatorname{trdeg}_kK(X)
=
\operatorname{trdeg}_kK(Y)
+
\operatorname{trdeg}_{K(Y)}K(X).
\]
Applying VI/30's dimension theorem,
\[
\dim X
=
\dim Y
+
\operatorname{trdeg}_{K(Y)}K(X).
\]
Hence
\[
\boxed{
\dim X-\dim Y
=
\operatorname{trdeg}_{K(Y)}K(X).
}
\]

The right-hand side measures the number of algebraically independent parameters remaining after the generic
point of \(Y\) has been fixed. This is why it is interpreted as the generic relative dimension of the dominant
map. In particular, the map is generically finite exactly when the relative transcendence degree is \(0\), and
then
\[
\dim X=\dim Y.
\]
\end{solution}

% VI30-DETAILED-CHALLENGE-03
\subsection*{Challenge VI.30.C3 solution}
\begin{solution}
Let
\[
X=\coprod_{n\ge0}\mathbb A^n_k.
\]
Its irreducible components are exactly the open-and-closed pieces \(\mathbb A^n_k\), and their dimensions are
unbounded. Therefore
\[
\dim X
=
\sup_{n\ge0}\dim\mathbb A^n_k
=
\sup_{n\ge0}n
=
\boxed{\infty}.
\]

Now let \(U\subseteq X\) be quasi-compact and open. Since
\[
U=\bigcup_{n\ge0}\left(U\cap\mathbb A^n_k\right)
\]
and the intersections are open subsets of \(U\), quasi-compactness gives a finite subcover. Because the components
are disjoint, this means that \(U\) meets only finitely many components, say those indexed by
\[
n_1,\ldots,n_r.
\]
Hence
\[
U=
\coprod_{j=1}^r\left(U\cap\mathbb A^{n_j}_k\right).
\]
Each piece has dimension at most \(n_j\), so
\[
\dim U
\le
\max_j n_j
<
\infty.
\]

Thus
\[
\boxed{\dim X=\infty}
\]
while every quasi-compact open subset has finite dimension. The example shows that global infinite dimension can
come from an unbounded family of finite-dimensional components rather than from an infinite chain living in one
quasi-compact region.
\end{solution}

% VI30-DETAILED-CHALLENGE-04
\subsection*{Challenge VI.30.C4 solution}
\begin{solution}
The product of integral \(k\)-schemes need not be integral when the ground field is not algebraically closed.

Take
\[
k=\mathbb R,
\qquad
X=Y=\Spec\mathbb C.
\]
Both are integral finite-type \(\mathbb R\)-schemes. But
\[
X\times_{\mathbb R}Y
=
\Spec(\mathbb C\otimes_{\mathbb R}\mathbb C),
\]
and
\[
\mathbb C\otimes_{\mathbb R}\mathbb C
\cong
\mathbb C\times\mathbb C.
\]
Hence the product is reducible.

A positive-dimensional version is
\[
X=Y=\Spec\mathbb C[t]
\]
viewed as \(\mathbb R\)-varieties. Then
\[
\mathbb C[t]\otimes_{\mathbb R}\mathbb C[s]
\cong
(\mathbb C\otimes_{\mathbb R}\mathbb C)[t,s]
\cong
\mathbb C[t,s]\times\mathbb C[t,s].
\]
Thus the product has two irreducible components, each of dimension \(2\), while
\[
\dim X=\dim Y=1.
\]

Accordingly, when the product is reducible one should not call the whole product a variety. The safe
dimension statement is componentwise: in the finite-type situation, the irreducible components of the product
have the expected dimension
\[
\dim X+\dim Y,
\]
under the standard hypotheses guaranteeing the usual product-dimension theorem. In the explicit example above,
each component has dimension
\[
2=1+1.
\]

If one assumes geometric integrality of one factor (and the usual finite-type hypotheses), the tensor product
remains a domain after the relevant scalar extension, eliminating this splitting issue.
\end{solution}

% VI30-DETAILED-CHALLENGE-05
\subsection*{Challenge VI.30.C5 solution}
\begin{solution}
The five statements have different hypothesis profiles.

\textbf{Open subsets cannot increase dimension.}
For every scheme and every open \(U\subseteq X\),
\[
\dim U\le\dim X.
\]
No finiteness hypothesis is needed. Take closures of chains of irreducible closed subsets.

\textbf{Every nonempty open preserves dimension.}
This is false for arbitrary schemes, false for reducible finite-type schemes, and even false for integral
Noetherian schemes. It is true for varieties:
\[
X\text{ integral finite type over }k
\quad\Longrightarrow\quad
\dim U=\dim X
\]
for every nonempty open \(U\).

\textbf{Proper closed subsets lower dimension.}
This fails for reducible schemes: one maximal-dimensional irreducible component can be a proper closed subset of
the whole scheme with equal dimension. If \(X\) is irreducible of finite dimension and
\(Y\subsetneq X\) is closed, then every chain in \(Y\) can be extended by \(X\), so
\[
\dim Y<\dim X.
\]

\textbf{Normalization preserves dimension.}
For an integral scheme with its standard normalization,
\[
\dim X^\nu=\dim X.
\]
Affinely this is the theorem that integral extensions preserve Krull dimension; the global statement follows by
affine covers.

\textbf{Finite surjective maps preserve dimension.}
For finite surjective morphisms in the affine case,
\[
\Spec B\to\Spec A,
\]
the ring \(B\) is integral over \(A/\ker(A\to B)\), and surjectivity makes the kernel invisible on the spectrum.
Thus
\[
\dim B=\dim A.
\]
The same statement globalizes over affine target covers.

\textbf{Dominant maps do not lower source dimension.}
For dominant morphisms of varieties,
\[
f:X\to Y,
\]
the field embedding
\[
K(Y)\hookrightarrow K(X)
\]
gives
\[
\dim X\ge\dim Y.
\]
The function-field proof requires the integral finite-type setting used to identify dimension with transcendence
degree.

The hypotheses are not uniform. Some statements are topological, while others come from integral extensions.
The strongest open-set and dominant-map formulas are specific to finite-type varieties.
\end{solution}
